\chapter{Alcohol Dependence (Alcoholism)}

\section*{Definition and Overview}
Alcohol dependence, commonly called alcoholism, is a chronic condition in which a person has lost control over their drinking. It is defined by a strong craving for alcohol, difficulty limiting how much is drunk, tolerance (needing more alcohol for the same effect), and withdrawal symptoms when drinking is reduced or stopped.

People with alcohol dependence often continue drinking despite serious harm to their health, work, or relationships. It is a recognised medical condition with a clear biological basis, not a moral failing, and effective treatments are available. Early recognition and treatment substantially improve the chance of recovery.

\section*{Epidemiology}
Alcohol dependence affects roughly 1 in 20 to 1 in 30 adults at some point in their lives and is about twice as common in men as in women. It most often begins in the late teens to the 30s.

People who start drinking young, binge drink, or have a family history of alcoholism are at higher risk. Alcohol dependence is a major cause of disability, hospital admission, and premature death, and it is frequently linked with depression, anxiety, and other substance use disorders.

\section*{Aetiology and Pathophysiology}
Alcohol dependence arises from an interaction of biological, psychological, and environmental factors.

\begin{itemize}
    \item \textbf{Genetics:} a family history of alcoholism is a strong risk factor, and heritability is thought to be around 50 percent
    \item \textbf{Brain chemistry:} repeated heavy drinking changes the brain's reward and stress systems, making alcohol increasingly necessary to feel normal
    \item \textbf{Tolerance:} the brain adapts to alcohol, so progressively more is needed to achieve the same effect
    \item \textbf{Withdrawal:} stopping drinking triggers unpleasant symptoms, which alcohol relieves, reinforcing the drinking cycle
    \item \textbf{Psychological factors:} drinking to cope with stress, low mood, or anxiety
    \item \textbf{Social factors:} heavy-drinking cultures, easy availability of alcohol, and peer groups that drink heavily
    \item \textbf{Early adversity:} childhood trauma, abuse, or neglect increases the risk of later dependence
\end{itemize}

\section*{Clinical Features}
\begin{itemize}
    \item A strong urge or craving to drink, sometimes present throughout the day
    \item Difficulty controlling when drinking starts and how much is drunk once it has
    \item Tolerance -- needing more alcohol to feel the same effect
    \item Withdrawal symptoms when drinking stops or is reduced: tremor, sweating, nausea, anxiety, and poor sleep
    \item Drinking to avoid withdrawal, including drinking first thing in the morning
    \item Spending a large amount of time obtaining, drinking, or recovering from alcohol
    \item Giving up other activities and responsibilities in order to drink
    \item Continuing to drink despite clear harm to health, work, or relationships
    \item Blackouts and memory lapses after drinking
\end{itemize}

\section*{Differential Diagnosis}
\begin{itemize}
    \item Hazardous or harmful drinking, in which drinking is risky but dependence is not yet established
    \item Binge drinking, which involves occasional heavy episodes rather than daily dependence
    \item Depression and anxiety, which can coexist with heavy drinking and cause it
    \item Dependence on other substances, such as benzodiazepines or opioids
    \item Medical conditions that mimic withdrawal symptoms, such as thyroid disease or low blood sugar
\end{itemize}

\section*{When to Seek Medical Care}
Seek help if drinking is hard to control, is causing health or relationship problems, or if withdrawal symptoms make it difficult to stop. Detoxifying from severe alcohol dependence at home can be dangerous and should be done under medical supervision.

\textbf{Call 000 (or local emergency services) for:}
\begin{itemize}
    \item A seizure or convulsion during withdrawal
    \item Severe confusion, hallucinations, or agitation (delirium tremens), which is a medical emergency
    \item Thoughts of suicide or self-harm
    \item Signs of alcohol poisoning, such as unresponsiveness, vomiting, or slow breathing
\end{itemize}

\section*{Diagnosis and Investigations}
\begin{itemize}
    \item \textbf{Clinical interview:} a detailed discussion of drinking patterns, craving, tolerance, withdrawal, and the impact on life
    \item \textbf{Screening tools:} questionnaires such as the AUDIT and CAGE
    \item \textbf{Physical examination:} looking for signs of liver disease, nerve damage, and malnutrition
    \item \textbf{Blood tests:} liver function tests, GGT, MCV, and carbohydrate-deficient transferrin can support the diagnosis
    \item \textbf{Mental health assessment:} evaluation for depression, anxiety, and the risk of harm
\end{itemize}

\section*{Management}
\begin{itemize}
    \item \textbf{Detoxification (withdrawal management):} supervised medical care, sometimes with medication, to withdraw from alcohol safely
    \item \textbf{Psychological therapies:} cognitive behaviour therapy and motivational interviewing to build motivation and prevent relapse
    \item \textbf{Medication:} naltrexone, acamprosate, and other medicines prescribed by a doctor can reduce craving or make drinking less rewarding
    \item \textbf{Mutual support groups:} such as Alcoholics Anonymous and SMART Recovery
    \item \textbf{Family involvement:} education and support for partners and family members
    \item \textbf{Treating co-existing problems:} such as depression, anxiety, or other substance use
    \item \textbf{Relapse prevention:} identifying triggers, building coping skills, and planning how to handle lapses
\end{itemize}

\section*{Complications}
\begin{itemize}
    \item Withdrawal complications, including seizures and delirium tremens
    \item Liver disease, including fatty liver, hepatitis, and cirrhosis
    \item Pancreatitis, gastritis, and other digestive problems
    \item Nerve damage and cognitive problems, including Wernicke--Korsakoff syndrome
    \item Heart disease, high blood pressure, stroke, and several cancers
    \item Injuries, accidents, and suicide
    \item Family breakdown, job loss, and financial problems
\end{itemize}

\section*{Prognosis and Outcomes}
Alcohol dependence is a chronic, relapsing condition, and relapse is common and does not mean that treatment has failed. Many people achieve sustained recovery, often after several attempts. Detoxification alone has a high relapse rate, but combining it with ongoing therapy, medication, and support substantially improves the chances of long-term recovery. Early treatment of the dependence and its complications improves the outlook.

\section*{Prevention and Self-Care}
\begin{itemize}
    \item \textbf{Low-risk drinking:} stay within low-risk guidelines and include alcohol-free days
    \item \textbf{Avoid binge drinking:} avoid heavy drinking episodes, and delay the first use of alcohol in young people
    \item \textbf{Coping:} address stress, low mood, and anxiety early rather than using alcohol to cope
    \item \textbf{Early help:} seek help as soon as you notice loss of control over drinking
    \item \textbf{Support:} ask family and friends for support, and involve them in recovery
    \item \textbf{Follow-up:} attend appointments and keep engaging with treatment after detoxification
    \item \textbf{Carers:} if you care for someone with alcohol dependence, encourage treatment and look after your own wellbeing
\end{itemize}

\section*{Sources and Further Reading}
The following sources provide reliable, up-to-date information on alcohol dependence.

\begin{itemize}
    \item NHS. \textit{Alcohol dependence: support and treatment.} https://www.nhs.uk
    \item World Health Organization. \textit{Alcohol use disorders and treatment guidance.} https://www.who.int
    \item Mayo Clinic. \textit{Alcohol use disorder: symptoms, causes, and treatment.} https://www.mayoclinic.org
    \item National Institute on Alcohol Abuse and Alcoholism. \textit{Understanding alcohol use disorder.} https://www.niaaa.nih.gov
    \item MedlinePlus. \textit{Alcohol use disorder.} https://medlineplus.gov
    \item Centers for Disease Control and Prevention. \textit{Alcohol and public health resources.} https://www.cdc.gov
\end{itemize}
